\scanpage{27}\mspage{IV}
2 - \uline{Rappel sur les sous-ensembles convexes de $\mathbb{R}^{2}$}.

\margin{!}\uline{Lemme II}. Si $A$ est une partie \inserted{\uline{finie}} de $\mathbb{R}^{2}$ contenant trois points affinement independants, il existe un polygone et un seul $\Gamma$ tel qu'il n'y ait pas trois sommets consecutifs colineaires \inserted{de $\Gamma$}, et \deleted{de façon} que l'ensemble sous-jacent à $\Gamma$ soit égal a $E(A)-\overset{\circ}{E(A)}$ ($E(A)=$ enveloppe convexe de $A$).

On a aussi $S(\Gamma)\subset A$.

\margin{jargon\\\tikz[x=.12cm,y=.12cm,baseline=0pt]{\draw (0,0)--(.12,-.10)--(.85,3.3);}\quad ?\\$\cdot$}%
% slanted
\uline{Demonstration} (\unsure{Esquisse|Esquissè|Esquissé}) Cela résulte du fait que $E(A)=\tikz[baseline=(q.base)]{\node[inner sep=1pt] (q) {$\cap$};\draw[rotate around={-12:(q.center)}] (q.center) ellipse (.17cm and .35cm);}$ demi espaces fermés correspondants à un ensemble de droites uniquem\uline{ent determinées par} la donnée de $A$, On remarque aussi que $S(\Gamma)=$ \inserted{ens. des} points extremaux de $E(A)$.

\margin{!!!\\partie \unsure{polyédr.l|polyédral|polyédrale} conv.}%
% slanted
\uline{Convention} : Le mot ensemble convexe désigne \uline{pour nous} l'enveloppe convexe d'un ensemble fini de points dans $\mathbb{R}^{2}$. \inserted{Soit $\mathbb{C}$ un \uline{ensemble convexe},} Si l'interieur est non-vide on munit $\partial\mathbb{C}=\mathbb{C}-\overset{\circ}{\mathbb{C}}$ de la structure de polygone décrite dans le lemme II. Il est clair qu'un \uline{ensemble convexe} $\mathbb{C}$ est compact et \inserted{par conséquent} $\partial\mathbb{C}=\mathbb{C}-\overset{\circ}{\mathbb{C}}$ aussi.
\margin{\unsure{intervertir|intervention} \unsure{avec|en} lemme II ! C'est une salade\ldots}

\uline{Lemme III}. Si $s,s'\in\partial\mathbb{C}$, $\overset{\circ}{\mathbb{C}}\ne\phi$, sont deux sommets consecutifs, la droite \deleted{$d$ qui} \inserted{$d$ qu'ils} déterminent partage $\mathbb{R}^{2}$ en deux demi-espaces fermés dont l'un contient $\overset{\circ}{\mathbb{C}}$.

\margin{?\quad ?\\vis\\1}%
\uline{Demonstration} : Cela est trivial quand on exprime $\mathbb{C}$ comme intersection de demi-espaces, si l'on tient compte que $\overset{\circ}{\mathbb{C}}$ est l'intersection des demi espaces ouverts correspondants, et qu'une arète de $\partial\mathbb{C}$ est toujours contenue dans un \uline{tel} demi-espace, c.q.f.d.

\uline{Lemme IV}. Soit $\mathbb{C}$ un \uline{ensemble convexe} à interieur non vide. Si $x\notin\mathbb{C}$ il existe au moins un point $y\in\mathbb{C}-S(\partial\mathbb{C})$ tel que $y$ est $(\mathbb{R}^{2}-C)$-visible de $x$.

\uline{Démonstration} : Choisissons $z\in\overset{\circ}{\mathbb{C}}$ tel que la droite determinée par $x$ et par $z$ \deleted{coupe $\partial\mathbb{C}$} ne contienne aucun sommet de $\partial\mathbb{C}$ (cela est possible car il n'y a qu'un nombre fini de Sommets). Comme $d\cap\mathbb{C}$ est un convexe fermé de $d$ à interieur \inserted{(relatif)} non-vide il est de la forme $d\cap C=[z_{1},z_{2}]$, $z_{1}\ne z_{2}$, il suffit de prendre $y$ \deleted{$\unsure{=y_{1}|=z_{1}}$} egal au point de $[z_{1},z_{2}]$ qui realise le minimum de la distence de $x$ a $[z_{1},z_{2}]$, c.q-f-d.
\margin{!}
% archive sheet number: 26
